To simulate the quasi-static electric potential distribution in the phantom head geometry,
we generated a three-dimensional finite element mesh using Gmsh~\cite{gmsh}. The mesh was
designed to accurately resolve regions with steep voltage gradients near the internal
electrodes and to capture the layered structure of the 3.2~mm conductive PLA shell,
while maintaining computational efficiency in the homogeneous gel interior. 
    
\paragraph*{Element Order and Type} We employed second-order tetrahedral elements (P2)
throughout the computational domain. Second-order elements are standard in bioelectric
field modeling because their quadratic interpolation yields substantially higher accuracy
than linear (P1) elements for comparable mesh sizes~\cite{zhang-2004}.
In our case, P2 elements allow accurate representation of smooth potential fields and 
bending effects with fewer elements. Studies, such as~\cite{zhang-2004}, showed that
second-order FEM provides “substantially enhanced numerical accuracy and computational
efficiency” relative to first-order FEM for similar meshes.
This property is particularly valuable for our 3.2~mm thick shell: using P2 elements,
a single layer of tetrahedra through the thickness suffices to capture the voltage drop
across the shell without excessive mesh refinement.
The mesh was configured with \texttt{Mesh.ElementOrder = 2} and \texttt{Mesh.HighOrderOptimize = 1}.
In Gmsh, \texttt{HighOrderOptimize=1} applies basic optimization of high-order node placement. 
A more aggressive setting (\texttt{=2}) would invoke additional elastic smoothing, but this requires 
significantly more memory; we therefore chose \texttt{=1} given our 16~GB RAM constraint.
We also set \texttt{Mesh.SecondOrderLinear = 1} so that mid-edge nodes on internal edges are placed by
simple linear interpolation. This enforces straight edges for the quadratic elements. Given that 
our phantom geometry includes primarily planar electrode faces and cylindrical guides, straight-edge P2
elements are sufficient and avoid unnecessary curvature. (Note that tetrahedral meshes inherently approximate
curved surfaces well~\cite{erdbrugger-2023}.)

\paragraph*{Global Mesh Resolution}
The literature on EEG and transcranial electrical stimulation (tES) modeling indicates that a 2~mm
 element size achieves excellent accuracy~\cite{erdbrugger-2023}. In a comprehensive head model convergence study,
 it is reported that reducing the mesh resolution from 16~mm to 2~mm shrinks the mean relative difference measure~(RDM)
from about 10.54 to 0.18~\cite{erdbrugger-2023}. Guided by this, we set a maximum element size of 2.0~mm in the homogeneous
 gel interior, where potential gradients are smooth and coarser resolution is acceptable.
  
  
\paragraph*{Refinement Near Internal Electrodes}
Each of the nine internal electrodes is modeled as a cylindrical contact (3.5~mm diameter, 12~mm tall) that imposes
a Dirichlet condition ($u=V_i$) on the gel boundary. Such voltage sources produce localized gradients
in the surrounding gel, so we refined the mesh near them using a distance-based field. Specifically, we used a
threshold field yielding a minimum element size of 0.5~mm for points within 2.0~mm of any electrode surface,
increasing linearly to 2.0~mm size by 8.0~mm away. In practice, this means:
\begin{itemize}
\item Element size = 0.5~mm for distance $\leq 2.0$~mm from electrodes
\item Element size = 2.0~mm for distance $\geq 8.0$~mm from electrodes
\item Linear transition between 2.0 and 8.0~mm
\end{itemize}

This graded refinement ensures that the mesh conforms closely to the curved electrode surfaces without an
overwhelming number of elements far away. In preliminary testing, we verified that a 0.5~mm near-electrode
size was sufficient to converge the potential distribution (smaller sizes produced negligible changes).

\paragraph*{Shell Refinement and Thin-Layer Capture}
The conductive PLA shell in the head phantom is nominally 3.2~mm thick and about ten times more conductive
than the interior gel (we measured $\sigma_{\rm shell}\approx 4.1$~S/m vs.\ $\sigma_{\rm gel}\approx 0.33$~S/m).
This high contrast across a thin layer requires fine resolution to avoid numerical artifacts. 
Indeed, previous studies note that very thin compartments can suffer spurious “current leakage” if under-resolve~\cite{erdbrugger-2023}.
To mitigate this, we applied another distance-based refinement around the shell–gel interface. We set 
a 0.5~mm element size (yielding roughly 6 linear elements through 3.2~mm) for distances up to 3.0~mm from the
shell, transitioning to 2.0~mm by 8.0~mm away. This means:
\begin{itemize}
\item Element size = 0.5~mm for distance $\leq 3.0$~mm from shell surfaces
\item Element size increases to 2.0~mm by 8.0~mm distance
\end{itemize}
With quadratic tetrahedra, 0.5~mm yields about 8 elements (counting mid-edge nodes) across the shell,
which is more than sufficient to capture the potential drop and avoid leakage. For reference, the SimNIBS
mesh-generation guidelines~\cite{simnibs-custom-meshing} also recommend using roughly 0.5~mm resolution to
resolve thin structures, and other studies showed that second-order elements can achieve this accuracy
with fewer layers than first-order elements~\cite{zhang-2004}.

\paragraph*{Curvature-Based Refinement}
In addition to these explicit fields, we enabled curvature-adaptive meshing (\texttt{Mesh.MeshSizeFromCurvature = 25}).
This instructs Gmsh to place more nodes along regions of high geometric curvature.
Curvature-adaptive refinement is a common strategy in head modeling, since features like facial contours
and the occipital scalp surface are smoothly curved~\cite{erdbrugger-2023}. 
By setting 25 elements per $2\pi$ radians, we ensure that our mesh automatically tightens around any remaining
curved geometry, further improving accuracy. 


\paragraph*{Refinement Field Combination}
We combined the electrode- and shell-based fields using a minimum operator (\texttt{Field[…] = Min}),
so that at any point in space the smallest (finest) prescribed element size is applied. This automatically 
handles overlapping regions.

\paragraph*{Mesh Quality Control}
After initial mesh generation, we performed iterative smoothing (\texttt{Mesh.Smoothing = 10}) 
and activated both Gmsh’s built-in optimizer and the Netgen optimizer (\texttt{Mesh.Optimize = 1}
and \texttt{Mesh.OptimizeNetgen = 1}). 

The physical volume tags were defined for the gel and shell, and physical surface tags for the nine electrodes,
to facilitate assignment of material properties and boundary conditions in the solver. The exterior shell surface 
was left untagged, taking the natural Neumann (insulating) condition of the solver.


\subsubsection{Final Mesh Specifications}
The final mesh was generated using Gmsh~4.11 and saved in MSH 4.1 format for
compatibility with FEniCSx (dolfinx). Meshing was completed in approximately 5
minutes on a machine with 16~GB RAM. The mesh statistics, as given by the meshing software, are:

\begin{itemize}
  \item Total nodes: 6,539,752
  \item Total tetrahedral elements: 4,587,720 (second-order, 10-node)
  \item Surface triangles: 12,274 (used for electrode and shell boundary tagging)
  \item Average scaled Jacobian (SICN): 0.871
  \item Average shape quality (SIGE): 0.890
\end{itemize}

These values indicate that the final mesh maintains a high element quality,
with minimal distortion or skewness, while respecting the local refinement
needs near electrodes and shell interfaces. The configuration reflects a practical
compromise between numerical accuracy and computational resource constraints.